\section{Conclusion}
In this study, we conducted a systematic survey on 67 primary studies using ML/DL models for software security vulnerability detection. We collected the papers from different journals and conference venues including 25 conferences and 12 journals. Our review is established based on five major research questions and a set of sub-research questions. We devised the research questions in a comprehensive manner where to cover various dimensions of software vulnerability detection. Our analysis of primary studies indicated that there is a booming trend in the growth of using ML/DL models for software vulnerability detection. Our deep analysis of data sources of primary studies revealed that 65.7\% of studies use benchmark data for software vulnerability detection. We also find 6 broad categories of DL models along with 14 classic ML models used in software vulnerability detection. The categories of DL models are classified as recurrent models, graph models, attention models, convolutional models, general models, and transformer models. RNNs are by far the most popular DNNs in software vulnerability detection. Our analysis also finds that RNNs with LSTM cells are the most popular network architectures in recurrent models, accounting for 8 primary studies. In the category of graph models, GGNN is the most popular DL model used by 4 primary studies. Our results on vulnerability types reveal that the most frequent type of vulnerability covered in existing studies is Improper Control of a Resource Through its Lifetime - (CWE-664) accounting for 18 primary studies. In conclusion, we have identified a collection of ongoing challenges that necessitate further exploration in future studies involving the utilization of ML/DL models for software vulnerability detection.